% Mathématiques 6e — cours V3
% Source LaTeX rendue côté Astro/KaTeX.

\section{Objectifs}

Dans ce chapitre, on apprend à :

\begin{itemize}
\item nommer et noter un angle ;
\item reconnaître les principaux types d’angles ;
\item mesurer et construire des angles avec un rapporteur ;
\item comprendre angles adjacents, opposés par le sommet et supplémentaires ;
\item construire et utiliser une bissectrice ;
\item construire des triangles à partir de données ;
\item utiliser les propriétés des triangles rectangle, isocèle et équilatéral ;
\item utiliser la somme des angles d’un triangle ;
\item construire le cercle circonscrit à un triangle ;
\item définir et construire une symétrie axiale.
\end{itemize}

\section{1. Un angle}

Un angle est formé par deux demi-droites ayant la même origine.

Si les demi-droites sont :

\[
[OA)
\]

et :

\[
[OB),
\]

l’angle se note :

\[
\widehat{AOB}.
\]

La lettre du sommet est placée au milieu de la notation.

\section{2. Mesurer un angle}

La mesure d’un angle s’exprime en degrés.

On note par exemple :

\[
\widehat{AOB}=65^\circ.
\]

Le rapporteur permet de mesurer un angle.

Pour l’utiliser correctement :

\begin{enumerate}
\item placer le centre du rapporteur sur le sommet de l’angle ;
\item aligner la graduation de départ avec un côté de l’angle ;
\item lire la graduation rencontrée par l’autre côté ;
\item choisir la bonne échelle du rapporteur.
\end{enumerate}

\section{3. Les principaux types d’angles}

Un angle nul mesure :

\[
0^\circ.
\]

Un angle aigu vérifie :

\[
0^\circ<\alpha<90^\circ.
\]

Un angle droit mesure :

\[
90^\circ.
\]

Un angle obtus vérifie :

\[
90^\circ<\alpha<180^\circ.
\]

Un angle plat mesure :

\[
180^\circ.
\]

Un angle plein mesure :

\[
360^\circ.
\]

Classer un angle revient donc à comparer sa mesure à des valeurs de référence.

\section{4. Exemple développé 1 — classer un angle}

On mesure :

\[
\alpha=128^\circ.
\]

Comme :

\[
90^\circ<128^\circ<180^\circ,
\]

l’angle est obtus.

Pour :

\[
\beta=35^\circ,
\]

on a :

\[
0^\circ<35^\circ<90^\circ.
\]

Donc $\beta$ est aigu.

\section{5. Angles adjacents}

Deux angles sont adjacents lorsqu’ils :

\begin{itemize}
\item ont le même sommet ;
\item ont un côté commun ;
\item sont situés de part et d’autre de ce côté commun sans se chevaucher.
\end{itemize}

Si :

\[
\widehat{AOC}=40^\circ
\]

et :

\[
\widehat{COB}=25^\circ,
\]

alors :

\[
\widehat{AOB}=40^\circ+25^\circ=65^\circ
\]

si les deux angles sont adjacents et forment ensemble $\widehat{AOB}$.

\section{6. Angles supplémentaires}

Deux angles sont supplémentaires lorsque la somme de leurs mesures vaut :

\[
180^\circ.
\]

Ainsi, si :

\[
\alpha=72^\circ,
\]

l’angle supplémentaire mesure :

\[
180^\circ-72^\circ=108^\circ.
\]

\section{7. Angles opposés par le sommet}

Lorsque deux droites se coupent, elles forment des angles opposés par le sommet.

Deux angles opposés par le sommet ont la même mesure.

Si un angle mesure :

\[
53^\circ,
\]

l’angle qui lui est opposé par le sommet mesure également :

\[
53^\circ.
\]

Les deux angles adjacents sont alors supplémentaires.

Ils mesurent :

\[
180^\circ-53^\circ=127^\circ.
\]

\section{8. Construire un angle}

Pour construire un angle de mesure donnée :

\begin{enumerate}
\item tracer une première demi-droite ;
\item placer le centre du rapporteur sur son origine ;
\item repérer la graduation correspondant à la mesure demandée ;
\item placer un point ;
\item tracer la seconde demi-droite passant par ce point.
\end{enumerate}

Une construction doit être suffisamment précise pour que la mesure soit retrouvée au rapporteur.

\section{9. Bissectrice d’un angle}

La bissectrice d’un angle est la demi-droite qui partage cet angle en deux angles de même mesure.

\coursefigure{/images/mathematiques/6eme/angles-bissectrice.svg}{Angle partagé en deux parties égales par sa bissectrice.}{Si une demi-droite est la bissectrice de $\widehat{AOB}$, alors les deux angles qu’elle forme avec les côtés de l’angle ont la même mesure.}

Si :

\[
\widehat{AOB}=84^\circ,
\]

et si $[Ox)$ est la bissectrice, alors :

\[
\widehat{AOx}
=
\widehat{xOB}
=
\dfrac{84^\circ}{2}
=
42^\circ.
\]

\section{10. Construire une bissectrice au rapporteur}

Une première méthode consiste à mesurer l’angle puis à construire la moitié de sa mesure.

Pour :

\[
\widehat{AOB}=70^\circ,
\]

on calcule :

\[
\dfrac{70^\circ}{2}=35^\circ.
\]

On construit alors une demi-droite faisant :

\[
35^\circ
\]

avec chacun des deux côtés.

\section{11. Construire une bissectrice au compas}

On peut aussi construire une bissectrice sans mesurer l’angle :

\begin{enumerate}
\item tracer un arc de cercle centré au sommet qui coupe les deux côtés ;
\item conserver la même ouverture de compas ;
\item tracer deux arcs centrés sur les deux points obtenus ;
\item relier le sommet à l’intersection des deux arcs.
\end{enumerate}

Cette construction repose sur une situation d’équidistance.

\section{12. Construire un triangle}

Un triangle peut être construit lorsque l’on possède suffisamment de données.

On peut par exemple connaître :

\begin{itemize}
\item les trois longueurs de côtés ;
\item deux longueurs et l’angle compris ;
\item une longueur et deux angles.
\end{itemize}

Chaque donnée doit être utilisée avec l’instrument adapté :

\begin{itemize}
\item règle ;
\item compas ;
\item rapporteur.
\end{itemize}

\section{13. Construire un triangle à partir de trois côtés}

Supposons :

\[
AB=6\ \text{cm},
\]

\[
AC=5\ \text{cm},
\]

\[
BC=4\ \text{cm}.
\]

Méthode :

\begin{enumerate}
\item tracer $[AB]$ ;
\item tracer le cercle de centre $A$ et de rayon :
\end{enumerate}

\[
5\ \text{cm};
\]

\begin{enumerate}
\item tracer le cercle de centre $B$ et de rayon :
\end{enumerate}

\[
4\ \text{cm};
\]

\begin{enumerate}
\item choisir un point d’intersection $C$ ;
\item tracer $[AC]$ et $[BC]$.
\end{enumerate}

Les distances imposées sont ainsi respectées.

\section{14. Triangle rectangle}

Un triangle rectangle possède un angle droit.

Si le triangle $ABC$ est rectangle en $A$ :

\[
\widehat{BAC}=90^\circ.
\]

Les deux autres angles sont aigus et leur somme vaut :

\[
90^\circ.
\]

\section{15. Triangle isocèle}

Un triangle isocèle possède deux côtés de même longueur.

Si :

\[
AB=AC,
\]

le triangle est isocèle en $A$.

Les angles à la base ont alors la même mesure :

\[
\widehat{ABC}
=
\widehat{BCA}.
\]

Cette propriété permet de calculer des angles sans les mesurer.

\section{16. Triangle équilatéral}

Un triangle équilatéral possède trois côtés de même longueur.

Il possède également trois angles de même mesure.

Comme la somme des angles d’un triangle vaut :

\[
180^\circ,
\]

chaque angle d’un triangle équilatéral mesure :

\[
\dfrac{180^\circ}{3}=60^\circ.
\]

\section{17. Somme des angles d’un triangle}

Dans tout triangle :

\[
\boxed{
\widehat{A}
+
\widehat{B}
+
\widehat{C}
=
180^\circ
}.
\]

Cette propriété est valable quelle que soit la forme du triangle.

\section{18. Exemple développé 2 — angle manquant}

Dans un triangle, deux angles mesurent :

\[
52^\circ
\]

et :

\[
71^\circ.
\]

Le troisième angle mesure :

\[
180^\circ-52^\circ-71^\circ.
\]

Donc :

\[
180^\circ-123^\circ=57^\circ.
\]

Le troisième angle vaut :

\[
\boxed{57^\circ}.
\]

\section{19. Exemple développé 3 — triangle isocèle}

Le triangle $ABC$ est isocèle en $A$ et :

\[
\widehat{BAC}=40^\circ.
\]

Les deux angles à la base sont égaux.

La somme restante vaut :

\[
180^\circ-40^\circ=140^\circ.
\]

On partage en deux :

\[
\dfrac{140^\circ}{2}=70^\circ.
\]

Donc :

\[
\widehat{ABC}
=
\widehat{BCA}
=
70^\circ.
\]

\section{20. Vérifier qu’une construction est possible}

Les données d’une construction doivent être compatibles.

Par exemple, si un triangle possède deux angles :

\[
110^\circ
\]

et :

\[
80^\circ,
\]

leur somme vaut déjà :

\[
190^\circ.
\]

Un tel triangle est impossible puisque la somme totale doit être :

\[
180^\circ.
\]

Avant de tracer, un calcul peut donc éviter une construction impossible.

\section{21. Médiatrices d’un triangle}

Chaque côté d’un triangle possède une médiatrice.

Les trois médiatrices d’un triangle se coupent en un même point.

On dit qu’elles sont \textbf{concourantes}.

Si leur point d’intersection est $O$, alors :

\[
OA=OB=OC.
\]

Le point $O$ est donc équidistant des trois sommets.

\section{22. Cercle circonscrit}

\coursefigure{/images/mathematiques/6eme/triangle-cercle-symetrie.svg}{Triangle inscrit dans un cercle dont le centre est l’intersection des médiatrices, accompagné d’un exemple de deux points symétriques.}{L’intersection des médiatrices d’un triangle est équidistante des trois sommets ; elle est donc le centre du cercle passant par ces trois sommets.}

Le cercle de centre $O$ et de rayon :

\[
OA
\]

passe alors par :

\[
A,\quad B,\quad C.
\]

Il est appelé \textbf{cercle circonscrit} au triangle.

\section{23. Construire le cercle circonscrit}

Pour construire le cercle circonscrit à un triangle $ABC$ :

\begin{enumerate}
\item construire la médiatrice de $[AB]$ ;
\item construire la médiatrice de $[AC]$ ;
\item noter $O$ leur point d’intersection ;
\item tracer le cercle de centre $O$ passant par $A$.
\end{enumerate}

La troisième médiatrice doit également passer par $O$.

\section{24. Symétrie axiale}

Deux points $A$ et $A'$ sont symétriques par rapport à une droite $(d)$ lorsque $(d)$ est la médiatrice du segment $[AA']$.

Cela implique :

\begin{itemize}
\item $(d)$ est perpendiculaire à $(AA')$ ;
\item $(d)$ passe par le milieu de $[AA']$.
\end{itemize}

Un point situé sur l’axe reste inchangé par la symétrie.

\section{25. Construire le symétrique d’un point}

Pour construire $A'$, symétrique de $A$ par rapport à $(d)$ :

\begin{enumerate}
\item tracer la perpendiculaire à $(d)$ passant par $A$ ;
\item noter $H$ son intersection avec $(d)$ ;
\item placer $A'$ de l’autre côté de $(d)$ de sorte que :
\end{enumerate}

\[
HA'=HA.
\]

Alors $H$ est le milieu de $[AA']$ et $(d)$ est sa médiatrice.

\section{26. Propriétés de la symétrie axiale}

La symétrie axiale conserve :

\begin{itemize}
\item les longueurs ;
\item les mesures d’angles ;
\item l’alignement ;
\item les périmètres ;
\item les aires ;
\item la forme des figures.
\end{itemize}

Si $A'$ et $B'$ sont les images de $A$ et $B$, alors :

\[
A'B'=AB.
\]

\section{27. Construire l’image d’une figure}

Pour construire le symétrique d’un polygone :

\begin{enumerate}
\item construire le symétrique de chaque sommet ;
\item relier les points images dans le même ordre.
\end{enumerate}

Il faut construire les sommets géométriquement ; recopier « à l’œil » ne garantit pas la symétrie.

\section{28. Axes de symétrie}

Une droite est un axe de symétrie d’une figure lorsque la symétrie par rapport à cette droite laisse la figure inchangée.

Un rectangle possède :

\[
2
\]

axes de symétrie.

Un carré en possède :

\[
4.
\]

Un triangle isocèle non équilatéral en possède :

\[
1.
\]

Un cercle possède une infinité d’axes de symétrie : toutes les droites passant par son centre.

\section{29. Méthode — résoudre un problème d’angles}

\begin{enumerate}
\item identifier les angles connus ;
\item repérer les propriétés utilisables ;
\item écrire une égalité de mesures ;
\item calculer ;
\item vérifier que la mesure finale est compatible avec le type d’angle.
\end{enumerate}

Dans un triangle, penser à :

\[
180^\circ.
\]

Dans un angle plat, penser également à :

\[
180^\circ.
\]

Autour d’un point, un tour complet vaut :

\[
360^\circ.
\]

\section{30. Méthode — construire un triangle}

\begin{enumerate}
\item lister les données ;
\item choisir un premier côté ;
\item traduire chaque autre donnée par un cercle, un angle ou une distance ;
\item construire les intersections ;
\item vérifier les longueurs et angles obtenus.
\end{enumerate}

La construction doit être lisible et justifiée.

\section{31. Erreurs fréquentes}

\subsection{Lire la mauvaise graduation du rapporteur}

Avant de lire, vérifier où se trouve :

\[
0^\circ
\]

sur le côté de départ.

\subsection{Confondre bissectrice et médiatrice}

La bissectrice concerne un angle.

La médiatrice concerne un segment.

\subsection{Oublier la somme des angles d’un triangle}

Une construction produisant une somme différente de :

\[
180^\circ
\]

est impossible.

\subsection{Croire que les médiatrices se coupent uniquement dans certains triangles}

Elles sont concourantes dans tout triangle.

\subsection{Construire une symétrie à l’œil}

La distance à l’axe doit être conservée de part et d’autre.

\subsection{Confondre diagonale et axe de symétrie}

Une diagonale d’un rectangle quelconque n’est pas un axe de symétrie.

\section{32. À retenir}

Les mesures de référence sont :

\[
0^\circ,\quad
90^\circ,\quad
180^\circ,\quad
360^\circ.
\]

La bissectrice partage un angle en deux angles de même mesure.

Dans tout triangle :

\[
\boxed{
\widehat{A}+\widehat{B}+\widehat{C}=180^\circ
}.
\]

Dans un triangle isocèle, les angles à la base sont égaux.

Dans un triangle équilatéral, chaque angle mesure :

\[
60^\circ.
\]

Les médiatrices d’un triangle sont concourantes et leur intersection est le centre du cercle circonscrit.

Deux points sont symétriques par rapport à une droite lorsque cette droite est la médiatrice du segment qui les relie.

La symétrie axiale conserve les longueurs et les angles.
